\begin{figure}[ht]
\centering
\begin{tikzpicture}[>=Stealth, font=\small, node distance=1.35cm]
  \node[draw, rounded corners, fill=blue!8, minimum width=3.8cm, minimum height=0.95cm] (e1)
    {$E[\ell]\cong (\Z/\ell\Z)^2$};
  \node[draw, rounded corners, fill=blue!8, below left=1.6cm and 2.4cm of e1, minimum width=3.8cm, minimum height=0.95cm] (e2)
    {$E[\ell^2]\cong (\Z/\ell^2\Z)^2$};
  \node[draw, rounded corners, fill=blue!8, below right=1.6cm and 2.4cm of e1, minimum width=3.8cm, minimum height=0.95cm] (e3)
    {$E[\ell^3]\cong (\Z/\ell^3\Z)^2$};
  \node[draw, rounded corners, fill=green!10, below=2.7cm of e1, minimum width=4.6cm, minimum height=0.95cm] (tate)
    {$T_\ell(E)=\varprojlim E[\ell^n]\cong \Zl^2$};
  \node[draw, rounded corners, fill=orange!12, right=4.8cm of tate, minimum width=4.8cm, minimum height=0.95cm] (rep)
    {$\rho_{E,\ell}\colon G_\Q\to \GL_2(\Zl)$};

  \draw[->, very thick] (e2) -- node[above left] {$[\ell]$} (e1);
  \draw[->, very thick] (e3) -- node[above right] {$[\ell]$} (e1);
  \draw[->, very thick] (tate) -- node[above] {Galois 作用} (rep);
  \draw[dashed, ->, thick] (e1) -- node[left] {逆极限} (tate);
  \draw[dashed, ->, thick] (e2) -- (tate);
  \draw[dashed, ->, thick] (e3) -- (tate);
\end{tikzpicture}
\caption{椭圆曲线的 $\ell^n$-挠点形成一个逆系统，其逆极限就是 Tate 模 $T_\ell(E)$；$G_\Q$ 在每一层的作用彼此相容，从而得到 $\ell$-进表示。}
\label{fig:ch09-tate-tower}
\end{figure}
